\section*{Kapitel 9 --- Transportsömmen och registeråtkomsten}
\addcontentsline{toc}{section}{Kapitel 9 --- Transportsömmen}

\solution{ex:9:1}{Förutsäg}
\ans{a} \code{RegTxDataHi} har index 4. Skrivning: \code{CmdWrite | 4} = \code{0x84}. Läsning:
bara indexet, alltså \code{0x04}.
\ans{b} \code{83 AA BB 00 00} --- kommandobyten \code{CmdWrite | 3}, sedan värdet MSB först.
\ans{c} \code{0x12345678}. Den första byten, \code{0xFF}, är svaret på kommandobyten och kastas;
de fyra följande sätts ihop.
\ans{d} Du får \code{0xFF123456}: värdet skiftas en byte, den sista databyten faller bort, och
skräpbyten hamnar överst. Symptomet är registervärden som ser ut att vara $256$ gånger för stora
med en oförklarlig \code{0xFF} i toppen.

\solution{ex:9:2}{Teckenutvidgningen}
Uppgiften ber om en kodändring; frågorna besvaras här.
\ans{a} \code{value} blir \code{0xFFFFFFFF80000000} trunkerat till 32 bitar, alltså
\code{0x80000000} --- men vägen dit går genom odefinierat beteende. \code{transfer()} returnerar
en \code{std::uint8_t}, som befordras till \code{int}. \code{0x80} skiftat 24 steg åt vänster
sätter teckenbiten i en 32-bitars \code{int}, och att skifta in i teckenbiten är odefinierat i
C++17. I praktiken blir resultatet ett negativt tal som sedan konverteras till
\code{std::uint32_t} med ettor där det inte ska finnas några.
\ans{b} För varje byte där bit 7 är \emph{låg}, alltså alla värden under \code{0x80}. Koden
fungerar därför perfekt för små identifierare och små DLC-värden.
\ans{c} Drivrutinen hade fungerat för \code{id} upp till \code{0x7F} i den översta databyten och
gått sönder däröver. Konkret: \code{0x123} fungerar, \code{0xAABB} i datafältet gör det inte.
Vid bring-upen ser det ut som att vissa frames kommer fram och andra inte, beroende på
\emph{innehållet} --- vilket är ett av de mest förvirrande symptom som finns, eftersom man letar
efter ett fel i överföringen när felet ligger i en typkonvertering.

\solution{ex:9:3}{En transport att felsöka med}
Uppgiften ber om kod. Utskriften ska gå att jämföra rad för rad mot exemplet i \kapref{ch:spi}:
fem bytes mellan ett \code{select} och ett \code{deselect}, kommandobyten först, värdet MSB först.

Två saker är värda att bygga in från början: skriv ut \emph{både} den skickade och den mottagna
byten på varje rad, så att en läsning går att följa, och räkna transaktionerna, så att en
\code{send()} går att kontrollera mot de sex den ska kosta.

\solution{ex:9:4}{Lagergränsen}
\ans{a} Nej. \code{ByteTransport} flyttar bytes och vet ingenting om frames. Ett beroende till
\code{frame.h} vore det första steget mot att den vet vad den transporterar.
\ans{b} Nej. Då vore det inte längre en transport utan en registeråtkomst, och den hör hemma i
\code{Spi} --- som är det enda lager som ska känna till registerkartan.
\ans{c} Nej, och det här är den som låter mest rimlig. Argumentet för är att transaktionen ändå
alltid är fem bytes, så varför inte låta transporten sköta den? Argumentet emot är att
\code{SS}-ramningen och femBytesformatet tillhör \emph{protokollet}, inte överföringen. Läggs de
i transporten måste varje transportimplementation --- även den skriptade testdubblaren --- kunna
protokollet, och då finns ingen söm kvar: man kan inte längre byta ut hur bytes flyttas utan att
samtidigt byta ut vad de betyder.
\ans{d} Nej. Att testa en privat metod direkt är nästan alltid ett tecken på att den borde vara
publik, eller på att testet testar fel sak. Här testas den genom \code{send()} och
\code{receive()}, och genom att inspektera vilka bytes som kom ut --- vilket är precis vad som
ska vara sant.
